% !TeX program = pdflatex
\documentclass[11pt,a4paper]{article}
\usepackage[utf8]{inputenc}
\usepackage[T1]{fontenc}
\usepackage[brazilian]{babel}
\usepackage{amsmath,amssymb,amsthm}
\usepackage[margin=2.4cm]{geometry}
\usepackage{booktabs}
\usepackage{array}
\usepackage{enumitem}
\usepackage{xcolor}
\usepackage[hidelinks,breaklinks]{hyperref}
\usepackage[expansion=false]{microtype}

\sloppy
\emergencystretch=2em

\newtheorem{definicao}{Definição}
\newtheorem{teorema}{Teorema}
\newtheorem{proposicao}{Proposição}
\newtheorem{lema}{Lema}
\newtheorem{observacao}{Observação}
\newtheorem{corolario}{Corolário}
\newtheorem{questao}{Questão}

\newcommand{\code}[1]{\texttt{\small #1}}
\definecolor{cinza}{gray}{0.35}
\newcommand{\abs}[1]{\lvert #1\rvert}

\title{Suporte Coordenado Recursivo\\
por Aninhamento de Pares}
\author{}
\date{}

\begin{document}
\maketitle

\begin{abstract}
\noindent
Formalizamos uma família recursiva de suportes coordenados construída por aninhamento de pares sobre um alfabeto finito arbitrário.
A base estrutural é sempre um par; a recursão acrescenta símbolos à esquerda.
A construção separa explicitamente
\[
\begin{array}{rcl}
\text{alfabeto} &\longrightarrow& B,\\
\text{codificação dos símbolos} &\longrightarrow& d,\\
\text{estrutura recursiva} &\longrightarrow& X^{(n)},\\
\text{coordenada} &\longrightarrow& C^{(n)},\\
\text{representação} &\longrightarrow& (b_1,b_2,\dots,b_n).
\end{array}
\]
A forma aninhada
\[
\bigl(b_1,(b_2,(b_3,\dots,(b_{n-1},b_n)\dots))\bigr)
\]
corresponde exactamente à expansão posicional
\[
d(b_1)+q\,d(b_2)+\cdots+q^{n-1}\,d(b_n).
\]
O caso dos bytes (\(q=256\)) aparece apenas como instanciação particular.
\end{abstract}

% ═══════════════════════════════════════════════════════════════════════════
\section{Motivação}

Em \texttt{campos.tex} um suporte coordenado é um par \((X,C)\) em que \(C\colon X\to\{0,\dots,\lvert X\rvert-1\}\) é bijecção.

Aqui construímos o suporte de forma recursiva por aninhamento de pares, com um alfabeto finito arbitrário.
A base estrutural é o par; a recursão acrescenta um novo símbolo à esquerda.
Cada nível contém uma cópia canónica do nível anterior, determinada por uma inclusão \(\iota_n\).

O objectivo é alinhar a estrutura sintáctica do aninhamento com a expansão posicional da coordenada:
\[
\bigl(b_1,(b_2,\dots,(b_{n-1},b_n)\dots)\bigr)
\quad\longleftrightarrow\quad
\sum_{i=1}^{n} q^{i-1}\,d(b_i).
\]

% ═══════════════════════════════════════════════════════════════════════════
\section{Definição}

\begin{definicao}[Alfabeto e codificação]
\label{def:alfabeto}
Seja \(B\) um conjunto finito com
\[
q
:=
\lvert B\rvert
\ge 2.
\]
Fixe uma bijecção (codificação dos símbolos)
\[
d\colon B\overset{\sim}{\longrightarrow}\{0,1,\dots,q-1\}.
\]
\end{definicao}

\begin{observacao}[Separação alfabeto/codificação]
\label{obs:alfabeto-codificacao}
O alfabeto \(B\) é um conjunto de símbolos; a codificação \(d\) atribui a cada símbolo um valor numérico.
Estas duas escolhas são independentes da estrutura recursiva do suporte.
Casos particulares:
\begin{itemize}[leftmargin=1.5em]
\item binário: \(q=2\), \(B=\{0,1\}\), \(d=\mathrm{id}\);
\item hexadecimal: \(q=16\);
\item byte: \(q=256\), \(B=\mathbb Z/256\mathbb Z\), \(d=\mathrm{id}\).
\end{itemize}
Nenhum deles é privilegiado pela definição.
\end{observacao}

\begin{definicao}[Família recursiva de suportes por aninhamento]
\label{def:suporte-aninhado}
Definimos indutivamente o conjunto \(X^{(n)}\) e a coordenada \(C^{(n)}\) para \(n\ge 2\).

\begin{enumerate}[leftmargin=1.5em]
\item \textbf{Base} (\(n=2\)):
\[
X^{(2)}
:=
B\times B,
\qquad
C^{(2)}(b_1,b_2)
:=
d(b_1)
+
q\,d(b_2).
\]
Assim \((X^{(2)},C^{(2)})\) é um suporte coordenado.

\item \textbf{Passo} (\(n\ge 3\)):
\[
X^{(n)}
:=
B\times X^{(n-1)}.
\]
Um elemento de \(X^{(n)}\) escreve-se
\[
x
=
(b,\, \xi)
\quad\text{com}\quad
b\in B,\quad
\xi\in X^{(n-1)}.
\]
A coordenada é definida por
\[
C^{(n)}(b,\, \xi)
\;:=\;
d(b)
\;+\;
q\, C^{(n-1)}(\xi).
\]
\end{enumerate}

Para cada \(n\ge 2\), o par \(\bigl(X^{(n)},\, C^{(n)}\bigr)\) é um suporte coordenado.
\end{definicao}

\begin{observacao}[Notação de aninhamento]
\label{obs:aninhamento}
Um estado de comprimento \(n\) escreve-se
\[
x
=
\bigl(
b_1,\,
\bigl(
b_2,\,
\bigl(
b_3,\,
\dots,\,
(b_{n-1},\, b_n)
\dots
\bigr)
\bigr)
\bigr),
\]
com \(b_i\in B\).
Aqui \(b_1\) é o símbolo menos significativo (peso \(q^0\)) e \(b_n\) o mais significativo (peso \(q^{n-1}\)).
A coordenada correspondente é
\[
C^{(n)}(x)
=
d(b_1)
+
q\,d(b_2)
+
\cdots
+
q^{n-1}\,d(b_n).
\]
A estrutura sintáctica e a expansão posicional estão alinhadas: o aninhamento à direita corresponde exactamente aos pesos crescentes.
\end{observacao}

% ═══════════════════════════════════════════════════════════════════════════
\section{Inclusão canónica}

\begin{definicao}[Inclusão canónica]
\label{def:iota-n}
Para cada \(n\ge 2\), a inclusão canónica é
\[
\iota_n\colon X^{(n)}\hookrightarrow X^{(n+1)},
\qquad
\iota_n(\xi)
:=
(b_\ast,\, \xi),
\]
onde \(b_\ast = d^{-1}(0)\) é o (único) símbolo de código zero.
A inclusão é determinada unicamente pela codificação \(d\).
Isto é: multiplica-se por \(q\) os pesos anteriores e insere-se \(0\) na posição de peso \(q^0\).
\end{definicao}

\begin{proposicao}[Escalonamento da coordenada sob inclusão]
\label{prop:escalonamento}
Para todo \(n\ge 2\),
\[
C^{(n+1)}\circ\iota_n
=
q\cdot C^{(n)}.
\]
Explicitamente:
\[
C^{(n+1)}\bigl(\iota_n(\xi)\bigr)
=
C^{(n+1)}(b_\ast,\, \xi)
=
d(b_\ast)
+
q\, C^{(n)}(\xi)
=
0
+
q\, C^{(n)}(\xi)
=
q\, C^{(n)}(\xi).
\]
\end{proposicao}

\begin{proof}
Imediato da definição de \(C^{(n+1)}\) e de \(\iota_n\).
\end{proof}

\begin{observacao}[Escalonamento de pesos]
\label{obs:escalonamento}
Ao contrário de uma construção que acrescenta o novo símbolo no extremo de peso máximo (e preserva a coordenada exactamente), aqui o novo símbolo é introduzido no extremo de peso mínimo.
Consequentemente a coordenada do nível anterior é multiplicada por \(q\): todos os pesos existentes sobem uma unidade.
A inclusão é injectiva e determinada unicamente pela codificação \(d\) (via \(b_\ast=d^{-1}(0)\)); a relação entre coordenadas é
\[
C^{(n+1)}\circ\iota_n
=
q\cdot C^{(n)}.
\]
\end{observacao}

\begin{corolario}[Coerência da família]
\label{cor:coerencia}
Para todo \(k\ge 1\),
\[
C^{(n+k)}\circ\iota_{n+k-1}\circ\cdots\circ\iota_n
=
q^k\cdot C^{(n)}.
\]
\end{corolario}

% ═══════════════════════════════════════════════════════════════════════════
\section{Propriedades fundamentais}

\begin{proposicao}[Bijecção]
\label{prop:bijecao}
Para todo \(n\ge 2\),
\[
C^{(n)}\colon X^{(n)}\overset{\sim}{\longrightarrow}\{0,1,\dots,q^{n}-1\}
\]
é bijecção.
Em particular \(\lvert X^{(n)}\rvert = q^{n}\).
\end{proposicao}

\begin{proof}
Por indução em \(n\).

\begin{itemize}[leftmargin=1.5em]
\item Base (\(n=2\)): \(C^{(2)}(b_1,b_2)=d(b_1)+q\,d(b_2)\) é a representação posicional clássica em base \(q\) com dois dígitos; a imagem é \(\{0,\dots,q^2-1\}\) e a aplicação é bijectiva.

\item Passo: suponha \(C^{(n)}\) bijecção sobre \(\{0,\dots,q^{n}-1\}\).
Então a aplicação
\[
(b,\, \xi)
\mapsto
d(b)
+
q\, C^{(n)}(\xi)
\]
é a representação posicional clássica em base \(q\), com \(d(b)\) no dígito de peso \(q^0\) e \(C^{(n)}(\xi)\) a ocupar os dígitos de pesos \(q^1,\dots,q^{n}\).
Como \(d(b)\) percorre \(\{0,\dots,q-1\}\) e \(C^{(n)}(\xi)\) percorre um bloco completo de comprimento \(q^{n}\), a imagem é \(\{0,\dots,q^{n+1}-1\}\) e a aplicação é injectiva (e portanto bijectiva).
\end{itemize}
\end{proof}

\begin{proposicao}[Ordem induzida]
\label{prop:ordem}
A ordem induzida por \(C^{(n)}\),
\[
\xi\le_{C^{(n)}}\eta
\quad\iff\quad
C^{(n)}(\xi)\le C^{(n)}(\eta),
\]
é total.
Sob a inclusão,
\[
\xi\le_{C^{(n)}}\eta
\quad\Longrightarrow\quad
\iota_n(\xi)\le_{C^{(n+1)}}\iota_n(\eta),
\]
porque \(C^{(n+1)}\circ\iota_n=q\cdot C^{(n)}\) e a multiplicação por \(q>0\) preserva a ordem.
\end{proposicao}

\begin{proof}
A totalidade segue de \(C^{(n)}\) ser bijecção para um intervalo de inteiros.
A preservação da ordem sob \(\iota_n\) é consequência directa da Proposição~\ref{prop:escalonamento}.
\end{proof}

% ═══════════════════════════════════════════════════════════════════════════
\section{Camadas distintas}

\begin{observacao}[Separação estrutural]
\label{obs:camadas}
A construção distingue explicitamente:
\begin{enumerate}[leftmargin=1.5em]
\item \textbf{Alfabeto} \(B\) --- conjunto finito de símbolos;
\item \textbf{Codificação} \(d\) --- bijecção que atribui valores numéricos aos símbolos;
\item \textbf{Estrutura recursiva} \(X^{(n)}\) --- o suporte propriamente dito, construído por aninhamento de pares a partir de um par-base;
\item \textbf{Coordenada} \(C^{(n)}\) --- a bijecção que atribui a cada estado um inteiro em \(\{0,\dots,q^{n}-1\}\);
\item \textbf{Representação} \((b_1,b_2,\dots,b_n)\) --- a sequência de símbolos que aparece quando se desaninha o estado.
\end{enumerate}
O suporte não é identificado com a representação numérica.
A coordenada é uma aplicação \emph{sobre} o suporte e depende da codificação \(d\).
A base estrutural é o par \(X^{(2)}=B\times B\); não existe nível de comprimento 1 na família.
\end{observacao}

\begin{observacao}[Comparação com o produto plano]
\label{obs:comparacao-produto}
Matematicamente, \(X^{(n)}\) é equipotente a \(B^{n}\) e transporta a mesma informação.
A diferença não está na cardinalidade nem na informação representada, mas na estrutura escolhida para o suporte:
em vez de apresentar cada nível directamente como um produto cartesiano plano \(B^{n}\), a construção o apresenta recursivamente por aninhamento de pares, com base no par.
Cada nível contém uma cópia canónica do nível anterior, determinada pela inclusão \(\iota_n\).
\end{observacao}

\begin{observacao}[Não se trata de um único objecto]
\label{obs:familia}
A construção produz uma \emph{família} recursiva de suportes
\[
X^{(2)},\ X^{(3)},\ X^{(4)},\ \dots
\]
na qual cada nível contém uma cópia canónica do nível anterior, determinada por \(\iota_n\).
Não se afirma a existência de um único suporte \(X^{(\infty)}\).
Uma eventual colagem indutiva
\[
X^{(\infty)}
:=
\varinjlim X^{(n)}
\]
seria uma estrutura adicional, desnecessária para a distinção entre suporte, coordenada e representação que se pretende aqui.
\end{observacao}

% ═══════════════════════════════════════════════════════════════════════════
\section{Espaço de campos sobre o suporte aninhado}

A construção do suporte aninhado introduz apenas estrutura combinatória e uma coordenada posicional.
Em particular, não se impõe estrutura linear sobre \(X^{(n)}\) nem sobre o alfabeto \(B\).
A estrutura linear pode ser introduzida posteriormente sobre os campos definidos no suporte.

\begin{definicao}[Espaço de campos escalares]
\label{def:espaco-campos}
Seja \(\mathbb K\) um corpo e seja \(X^{(n)}\) um suporte da família recursiva.
Definimos o espaço de campos escalares sobre \(X^{(n)}\) por
\[
\mathcal V_n
\;:=\;
\mathbb K^{X^{(n)}}
=
\{V\mid V\colon X^{(n)}\to\mathbb K\}.
\]
As operações são definidas ponto a ponto:
\[
(V+W)(x)
:=
V(x)+W(x),
\qquad
(\lambda V)(x)
:=
\lambda\,V(x),
\]
para \(V,W\in\mathcal V_n\), \(\lambda\in\mathbb K\) e \(x\in X^{(n)}\).
Com essas operações, \(\mathcal V_n\) é um espaço vectorial sobre \(\mathbb K\).
\end{definicao}

\begin{observacao}[Separação entre suporte e campo]
\label{obs:suporte-campo}
O suporte \(X^{(n)}\) e o espaço de campos \(\mathcal V_n\) desempenham papéis distintos:
\[
x\in X^{(n)}
\]
é um elemento do suporte, enquanto
\[
V\in\mathcal V_n
\]
é um campo definido sobre o suporte e
\[
V(x)\in\mathbb K
\]
é o valor desse campo em \(x\).
Assim
\[
\boxed{X^{(n)}\neq\mathcal V_n}
\]
e nenhuma estrutura linear é atribuída ao suporte apenas pela definição de \(\mathcal V_n\).
\end{observacao}

\begin{proposicao}[Dimensão]
\label{prop:dimensao-campos}
Para todo \(n\ge 2\),
\[
\dim_{\mathbb K}\mathcal V_n
=
\lvert X^{(n)}\rvert
=
q^{n}.
\]
\end{proposicao}

\begin{proof}
Como \(X^{(n)}\) é finito, o espaço de funções \(\mathcal V_n=\mathbb K^{X^{(n)}}\) possui uma base canónica indexada pelos elementos de \(X^{(n)}\).
Para cada \(x\in X^{(n)}\), defina
\[
e_x(y)
:=
\begin{cases}
1 & \text{se }y=x,\\
0 & \text{se }y\neq x.
\end{cases}
\]
Então
\[
\mathcal E_n
:=
\{e_x\mid x\in X^{(n)}\}
\]
é uma base de \(\mathcal V_n\).
Como \(\lvert X^{(n)}\rvert=q^{n}\), segue que
\[
\dim_{\mathbb K}\mathcal V_n
=
q^{n}.
\]
\end{proof}

\begin{observacao}[A coordenada como índice da base]
\label{obs:coordenada-base}
A coordenada
\[
C^{(n)}\colon X^{(n)}\overset{\sim}{\longrightarrow}\{0,\dots,q^{n}-1\}
\]
permite indexar a base canónica por inteiros.
Escrevendo
\[
e_c
:=
e_{(C^{(n)})^{-1}(c)},
\]
obtemos
\[
\mathcal E_n
=
\{e_c\mid 0\le c<q^{n}\}.
\]
Consequentemente, todo campo \(V\in\mathcal V_n\) possui a expansão
\[
V
=
\sum_{c=0}^{q^{n}-1}
V\bigl((C^{(n)})^{-1}(c)\bigr)\,e_c.
\]
A coordenada, portanto, indexa a base, mas não identifica o suporte com a base.
\end{observacao}

\begin{observacao}[Representação coordenada do campo]
\label{obs:campo-coordenado}
A bijecção \(C^{(n)}\) permite transportar um campo \(V\in\mathcal V_n\) para a representação coordenada
\[
\widehat V
\colon
\{0,\dots,q^{n}-1\}
\to
\mathbb K,
\]
definida por
\[
\widehat V(c)
:=
V\bigl((C^{(n)})^{-1}(c)\bigr).
\]
Assim,
\[
V
\longleftrightarrow
\widehat V
\]
é apenas uma mudança de índice determinada pela bijecção \(C^{(n)}\).
Não se introduz um novo campo nem uma nova grandeza: trata-se do mesmo campo escrito segundo a coordenada escolhida.
\end{observacao}

\begin{observacao}[Valores vectoriais]
\label{obs:valores-vectoriais}
A construção não depende de os valores dos campos serem escalares.
Se \(r\ge 1\), pode-se considerar
\[
\mathcal V_n^{(r)}
:=
(\mathbb K^r)^{X^{(n)}}.
\]
Um elemento \(V\in\mathcal V_n^{(r)}\) associa a cada \(x\in X^{(n)}\) um vector \(V(x)\in\mathbb K^r\).
Nesse caso,
\[
\dim_{\mathbb K}\mathcal V_n^{(r)}
=
r\cdot q^{n}.
\]
A estrutura do suporte e a coordenada permanecem inalteradas; somente o contradomínio dos campos é ampliado.
\end{observacao}

\begin{observacao}[Extensão entre níveis]
\label{obs:extensao-campo}
A inclusão canónica
\[
\iota_n\colon X^{(n)}\hookrightarrow X^{(n+1)}
\]
permite definir extensões de campos, mas tal extensão é uma construção adicional e não uma propriedade intrínseca do suporte.
Dado \(V\colon X^{(n)}\to\mathbb K\), podemos, por exemplo, escolher a extensão
\[
\widetilde V\colon X^{(n+1)}\to\mathbb K
\]
definida por
\[
\widetilde V(b,\xi)
:=
\begin{cases}
V(\xi) & \text{se }d(b)=0,\\
0 & \text{se }d(b)\neq 0.
\end{cases}
\]
Essa escolha preserva os valores do campo na cópia canónica \(\iota_n(X^{(n)})\subseteq X^{(n+1)}\) e atribui valor nulo às posições acrescentadas.
Nenhuma unicidade dessa extensão é afirmada.
\end{observacao}

\begin{observacao}[Camadas]
\label{obs:camadas-campo}
A construção pode ser resumida pelas seguintes camadas distintas:
\[
\begin{array}{rcl}
B &\longrightarrow& \text{alfabeto},\\[0.3em]
d &\longrightarrow& \text{codificação dos símbolos},\\[0.3em]
X^{(n)} &\longrightarrow& \text{suporte},\\[0.3em]
C^{(n)} &\longrightarrow& \text{coordenada},\\[0.3em]
\mathcal V_n=\mathbb K^{X^{(n)}} &\longrightarrow& \text{espaço de campos (vectorial)},\\[0.3em]
V(x) &\longrightarrow& \text{valor do campo}.
\end{array}
\]
A passagem de uma camada à seguinte não identifica os respectivos objectos.
Em particular, a estrutura linear pertence a \(\mathcal V_n\), e não é introduzida no suporte \(X^{(n)}\).
Quando o contradomínio é ele próprio um corpo, a camada de valores admite ainda uma estrutura de álgebra (secção seguinte).
\end{observacao}

% ═══════════════════════════════════════════════════════════════════════════
\section{Álgebra de campos sobre o suporte}

Quando o contradomínio dos campos é ele próprio um corpo \(Y\), a estrutura deixa de ser apenas vectorial: a multiplicação de \(Y\) induz uma multiplicação ponto a ponto no espaço de funções.

\begin{definicao}[Álgebra de campos]
\label{def:algebra-campos}
Seja \(Y\) um corpo e seja \(X^{(n)}\) um suporte da família recursiva.
Definimos
\[
\mathcal A_n(Y)
\;:=\;
Y^{X^{(n)}}
=
\{V\mid V\colon X^{(n)}\to Y\}.
\]
As operações são definidas ponto a ponto:
\begin{align*}
(V+W)(x)
&:=
V(x)+W(x),\\
(VW)(x)
&:=
V(x)\,W(x),\\
(\lambda V)(x)
&:=
\lambda\,V(x),\\
\mathbf 1(x)
&:=
1_Y,
\end{align*}
para \(V,W\in\mathcal A_n(Y)\), \(\lambda\in Y\) e \(x\in X^{(n)}\).
\end{definicao}

\begin{proposicao}[Base de idempotentes]
\label{prop:idempotentes}
Para cada \(x\in X^{(n)}\) defina
\[
e_x\colon X^{(n)}\to Y,
\qquad
e_x(y)
=
\begin{cases}
1_Y & \text{se }y=x,\\
0 & \text{se }y\neq x.
\end{cases}
\]
O conjunto \(\{e_x:x\in X^{(n)}\}\) é uma base de \(\mathcal A_n(Y)\) como espaço vectorial sobre \(Y\), e satisfaz
\[
e_x\,e_y
=
\begin{cases}
e_x & \text{se }x=y,\\
0 & \text{se }x\neq y,
\end{cases}
\]
bem como
\[
\mathbf 1
=
\sum_{x\in X^{(n)}} e_x.
\]
\end{proposicao}

\begin{proof}
Todo \(V\in\mathcal A_n(Y)\) escreve-se
\[
V
=
\sum_{x\in X^{(n)}}
V(x)\,e_x,
\]
pois, para cada \(y\in X^{(n)}\),
\[
\biggl(\sum_{x\in X^{(n)}}V(x)\,e_x\biggr)(y)
=
V(y).
\]
Se \(\sum_{x}\lambda_x e_x=0\) com \(\lambda_x\in Y\), então, avaliando em \(y\), obtém-se \(\lambda_y=0\). Logo \(\{e_x\}\) é linearmente independente e, portanto, base.
A relação de idempotência e ortogonalidade é verificação directa da definição ponto a ponto.
A decomposição da unidade é imediata.
\end{proof}

\begin{proposicao}[Estrutura de álgebra]
\label{prop:algebra}
\(\mathcal A_n(Y)\) é uma \(Y\)-álgebra comutativa unitária.
Além disso,
\[
\dim_Y\mathcal A_n(Y)
=
\lvert X^{(n)}\rvert
=
q^{n}.
\]
\end{proposicao}

\begin{proof}
As operações ponto a ponto herdam de \(Y\) a comutatividade, a associatividade e a distributividade.
A unidade é o campo constante \(\mathbf 1\).
A dimensão segue de \(\lvert X^{(n)}\rvert=q^{n}\) e da base de idempotentes \(\{e_x\}\) estabelecida na Proposição~\ref{prop:idempotentes}.
\end{proof}

\begin{observacao}[Produto de cópias]
\label{obs:produto-copias}
Como \(X^{(n)}\) é finito,
\[
\mathcal A_n(Y)
\;\cong\;
\prod_{x\in X^{(n)}} Y,
\]
com operações componente a componente.
A estrutura de álgebra é, portanto, a do produto finito de cópias de \(Y\).
\end{observacao}

\begin{observacao}[Separação de camadas]
\label{obs:separacao-algebra}
A álgebra \(\mathcal A_n(Y)\) vive no espaço de funções sobre o suporte.
Nenhuma estrutura algébrica é introduzida em \(B\) nem em \(X^{(n)}\).
O corpo de valores \(Y\) é uma escolha adicional, independente do alfabeto e da codificação \(d\).
Em particular, \(Y\) pode ser distinto de qualquer corpo usado noutras camadas do formalismo.
\end{observacao}

\begin{observacao}[Relação com o espaço vectorial]
\label{obs:relacao-Vn}
Quando \(Y=\mathbb K\), temos a identificação
\[
\mathcal A_n(\mathbb K)
=
\mathcal V_n
\]
como espaços vectoriais.
A multiplicação ponto a ponto é a estrutura adicional que só aparece quando se toma o contradomínio como corpo e se activa a operação \((VW)(x)=V(x)W(x)\).
\end{observacao}

% ═══════════════════════════════════════════════════════════════════════════
\section{Corpo de valores complexos por extensão quadrática}

A construção do corpo de valores é independente do suporte posicional \(X^{(n)}\).
Tratamos \(\mathbb R\) como dado e construímos o corpo complexo por quociente de polinómios.

\begin{definicao}[Anel de polinómios e ideal]
\label{def:RX}
Seja \(\mathbb R[X]\) o anel dos polinómios com coeficientes reais na indeterminada \(X\).
Considere o ideal principal gerado por \(X^2+1\):
\[
(X^2+1)
=
\{(X^2+1)\,g
\mid
g\in\mathbb R[X]\}.
\]
\end{definicao}

\begin{definicao}[Relação de congruência]
\label{def:congruencia}
Declare
\[
f\sim g
\quad\Longleftrightarrow\quad
f-g\in(X^2+1).
\]
Isto é uma relação de equivalência em \(\mathbb R[X]\).
\end{definicao}

\begin{definicao}[Quociente]
\label{def:Cconstr}
O corpo de valores complexos por construção é o anel quociente
\[
\mathbb C_{\mathrm{constr}}
\;:=\;
\mathbb R[X]/(X^2+1).
\]
Denote por \([f]\) a classe de equivalência de \(f\in\mathbb R[X]\).
\end{definicao}

\begin{proposicao}[Elemento \(i\)]
\label{prop:i}
Defina
\[
i
\;:=\;
[X].
\]
Então
\[
i^2
=
-1.
\]
\end{proposicao}

\begin{proof}
\[
i^2+1
=
[X]^2+[1]
=
[X^2+1]
=
[0],
\]
porque \(X^2+1\in(X^2+1)\). Logo \(i^2=-1\).
\end{proof}

\begin{proposicao}[Forma normal]
\label{prop:forma-normal}
Para todo \(f\in\mathbb R[X]\) existem únicos \(a,b\in\mathbb R\) tais que
\[
[f]
=
[a+bX].
\]
\end{proposicao}

\begin{proof}
Pelo algoritmo de divisão euclidiana em \(\mathbb R[X]\), existem únicos \(q,r\in\mathbb R[X]\) com
\[
f
=
(X^2+1)\,q
+
r
\quad\text{e}\quad
\deg r<2.
\]
Assim \(r=a+bX\) para únicos \(a,b\in\mathbb R\), e
\[
[f]
=
[r]
=
[a+bX].
\]
\end{proof}

\begin{observacao}[Identificação]
\label{obs:identificacao-C}
A forma normal permite a identificação
\[
[a+bX]
\;\longleftrightarrow\;
a+bi,
\]
onde \(i=[X]\). Esta é uma reescrita das classes, não uma definição prévia de \(\mathbb C\).
\end{observacao}

\begin{proposicao}[Multiplicação nas classes]
\label{prop:mult-classes}
Para \(a,b,c,d\in\mathbb R\),
\[
[a+bX]\,[c+dX]
=
[(ac-bd)+(ad+bc)X].
\]
\end{proposicao}

\begin{proof}
O produto dos representantes é
\[
(a+bX)(c+dX)
=
ac+(ad+bc)X+bd\,X^2.
\]
Como \(X^2\sim -1\), tem-se \(bd\,X^2\sim -bd\), logo
\[
[(a+bX)(c+dX)]
=
[(ac-bd)+(ad+bc)X].
\]
\end{proof}

\begin{observacao}[A fórmula familiar como leitura]
\label{obs:formula-leitura}
A expressão
\[
(a+bi)(c+di)
=
(ac-bd)+(ad+bc)i
\]
não é o ponto de partida: é o resultado da multiplicação de classes lida na forma normal.
\end{observacao}

\begin{proposicao}[\(\mathbb C_{\mathrm{constr}}\) é um corpo]
\label{prop:C-corpo}
\(\mathbb C_{\mathrm{constr}}\) é um corpo.
\end{proposicao}

\begin{proof}
\(\mathbb R[X]\) é um domínio de integridade euclidiano e \(X^2+1\) é irredutível sobre \(\mathbb R\) (não possui raízes reais). Logo o ideal \((X^2+1)\) é maximal e o quociente é um corpo.
\end{proof}

\begin{observacao}[Campos complexos sobre o suporte]
\label{obs:An-C}
Tomando \(Y=\mathbb C_{\mathrm{constr}}\) na álgebra de campos, obtém-se
\[
\mathcal A_n(\mathbb C_{\mathrm{constr}})
=
\mathbb C_{\mathrm{constr}}^{X^{(n)}}.
\]
A multiplicação de dois campos \(Z,W\) é
\[
(ZW)(x)
=
Z(x)\,W(x),
\]
onde o produto à direita é a multiplicação de classes já construída.
Quando se escolhe a forma normal em cada ponto,
\[
Z(x)=a_x+b_x i,
\qquad
W(x)=c_x+d_x i,
\]
recupera-se a fórmula clássica
\[
(ZW)(x)
=
(a_xc_x-b_xd_x)
+
(a_xd_x+b_xc_x)i.
\]
Nenhuma estrutura algébrica é introduzida no suporte \(X^{(n)}\).
\end{observacao}

\begin{observacao}[Independência das camadas]
\label{obs:independencia-C}
A construção
\[
\mathbb R
\;\longrightarrow\;
\mathbb R[X]
\;\longrightarrow\;
(X^2+1)
\;\longrightarrow\;
\mathbb R[X]/(X^2+1)
\;\longrightarrow\;
a+bi
\]
pertence exclusivamente à camada dos corpos de valores.
O suporte posicional \(X^{(n)}\) e a coordenada \(C^{(n)}\) permanecem inalterados.
A etapa \(\mathbb Q\to\mathbb R\) (cortes de Dedekind ou classes de Cauchy) pode ser documentada à parte; para o presente formalismo trata-se \(\mathbb R\) como dado.
\end{observacao}

% ═══════════════════════════════════════════════════════════════════════════
\section{Relação com \texttt{campos.tex}}

A família \(\bigl\{\bigl(X^{(n)},C^{(n)}\bigr)\bigr\}_{n\ge 2}\) satisfaz as condições de suporte coordenado de \texttt{campos.tex} (Definição de suporte coordenado).

A orientação do aninhamento (produto à esquerda, símbolo novo com peso mínimo) produz a relação
\[
C^{(n+1)}\circ\iota_n
=
q\cdot C^{(n)}.
\]
Isto corresponde a multiplicar por \(q\) os pesos anteriores e inserir \(0\) na posição de peso \(q^0\).

O caso dos bytes corresponde à instanciação
\[
q=256,
\qquad
B=\mathbb Z/256\mathbb Z,
\qquad
d=\mathrm{id}.
\]

A base estrutural é o par \(X^{(2)}=B\times B\); não há nível de comprimento unitário na família.
A forma aninhada e a expansão posicional ficam perfeitamente alinhadas:
\[
\bigl(b_1,(b_2,\dots,(b_{n-1},b_n)\dots)\bigr)
\quad\longleftrightarrow\quad
\sum_{i=1}^{n} q^{i-1}\,d(b_i).
\]

\end{document}
